% Migrated from days/day04.tex.  Paste your original problem, solution,
% remark and takeaway bodies into the slots below -- the wrappers are the only
% thing that changed.  Replace "the Day 4 result" style cross-references with
% \qref{Q004}, and recurring lemmas with \usetool{T1}.
\begin{entry}{Q004}{A pure marginal forces a product}

\entrymeta{
  topic   = {Bipartite states; supports; positivity},
  est     = $\sim$10 minutes,
  origin  = Original,
  source  = {Original problem},
  prereq  = {Partial trace},
  status  = written-up,
}

\begin{problem}
Let $\rho\in\Dop{\mathcal H_A\otimes\mathcal H_B}$, with $\mathcal H_A,\mathcal H_B$
finite dimensional, and suppose the reduced state $\rho_A=\Tr_B\rho$ is \emph{pure},
say $\rho_A=\proj{\psi_A}$. Prove that $\rho$ is a product state:
\[
  \rho \;=\; \proj{\psi_A}\otimes\rho_B,
  \qquad \rho_B=\Tr_A\rho .
\]
\end{problem}

\begin{solution}
For $p=\ket{\phi}\bra{\phi}$ a pure state, suppose it is non-separable.
We can look at its Schmidt decomposition, then 
$p= \sum_i \lambda_i \ket{i_A}\ket{i_B}$ with Schmidt rank at least 2, the partial trace over B is 
$p_A = \Tr_B p = \sum_i \lambda_i^2 \ket{i_A}\bra{i_A}$.

Since we have that $\sum_i \lambda_i^2 =1$ and the Schmidt rank is greater than 1, we have that $\lambda_i^2 < 1$  (for a pure state, the Schmidt rank is 1, so we have only a $\lambda_1=1$ which is not altered when squared)
So square each summand, $\lambda_i^2 \rightarrow \lambda_i^4$, will decrease the sum, so
$\Tr( p_A^2) = \sum_i \lambda_i^4 < 1$, meaning it is not a pure state.  
By contrapositive if $\rho$ is a pure state, and $\Tr_B(\rho)$ is  pure state, then is separable with $\rho =\ket{\psi_A}\otimes\ket{\psi_B}$.

Now we got look into $\rho$ mixed. Let $\ket{i_A}$ the a orthogonal basis on $\mathcal{H}_A$ with $\ket{0_A}=\ket{\psi}$. 
$\rho$ is diagonalizable, $\rho=\sum_i \rho_i \ket{\phi_i}\bra{\phi_i}$ with $\rho_i>0$ and $\sum_i \rho_i=1$.
We have that
\begin{equation}
    \ket{0_A}\bra{0_A}=\sum_i \rho_i \Tr_B(\ket{\phi_i}\bra{\phi_i})=\sum_i \rho_i x_i
    \label{check}
\end{equation}
where $x_i=\Tr_B(\ket{\phi_i}\bra{\phi_i}) \in \Dop{\mathcal H_A}$.

Thus $\sum_i \rho_i \bra{k_A}x_i\ket{k_A}=0$
for every $k>0$, but since each $\rho_i$ is positive, we need that $\bra{k_A}x_i\ket{k_A}=0$, so
$x_i$ must be proportional to $\ket{0_A}\bra{0_A}$ due to $x_i$ positivity, and actually equal due to normalization (it is a state). Which is consistent with \eqref{check} and $\sum_i \rho_i=1$.

But now this means that each pure state $\ket{\phi_i}$ has as partial trace a pure state $\ket{0_A}$, so we can use the result for pure states to guarantee that $\ket{\phi_i}=\ket{0_A} \otimes \ket{i_B}$
$$\rho_A= \ket{0_A}\bra{0_A} \otimes \sum_i \rho_i \ket{i_B}\bra{i_B}.$$

Notice $\ket{i_B}$ is not necessarily a basis of $\mathcal H_B$, it is just abuse of notation.
$X$ positive, implies $X=Y^\dagger Y$, so $\bra{\psi}X\ket{\psi}=||Y\ket{\psi}||^2=0$, implies $X\ket{\psi}=Y^\dagger Y\ket{\psi}=0$.
\end{solution}

\begin{takeaway}

\end{takeaway}

\attempt{
  date       = 2026-09-10,
  time       = 20 min,
  difficulty = 2,
  howitwent  = ,
}

\end{entry}
